\chapter{Your Toolkit and Environment Setup}
\label{ch:python}

\chapterguide%
  {Create folders and run basic \textbf{terminal commands}.}
  {organize data and notebooks, configure an environment, install dependencies, and run labs reproducibly.}

\section{The course project layout}

Use \texttt{UTP-"COURSE-NAME"/} as the project root, with these folders:

\begin{lstlisting}[style=utpterminal]
Lab01/
  student_records.cpp
  union_review.cpp
Lab02/
  matrix_multiply.cpp
  matrix_operations.cpp
Lab03/
  doubly_linked_list.cpp
  singly_linked_list.cpp
Lab04/
  circular_linked_list.cpp
  list_comparison.cpp
Lab05/
  array_stack.cpp
  infix_to_postfix.cpp
  linked_stack.cpp
Lab06/
  circular_queue.cpp
  linear_queue.cpp
  linked_queue.cpp
  priority_queue.cpp
Lab07/
  binary_tree.cpp
Lab08/
  bst.cpp
  heap.cpp
Lab09/
  graph_traversal.cpp
  minimum_spanning_tree.cpp
Lab10/
  binary_search.cpp
  heap_sort.cpp
  merge_three_arrays.cpp
  quick_sort_descending.cpp
Lab11/
  hashing_collisions.cpp
README.md
.gitignore
\end{lstlisting}
\clearpage
\section{C++ compiler setup}

Install the compiler for your operating system. The course examples use C++17; an editor or IDE is optional.

\medskip
\noindent\textbf{Windows}\par
\noindent Install \href{https://www.msys2.org/}{\textbf{MSYS2}}, open its
\textbf{UCRT64} terminal, then install and check GCC:

\begin{lstlisting}[style=utpterminal]
pacman -S mingw-w64-ucrt-x86_64-gcc
g++ --version
\end{lstlisting}

\noindent To run \texttt{g++} from PowerShell, Command Prompt, or an IDE
terminal, add the UCRT64 compiler folder to your Windows user \texttt{Path}:

\begin{enumerate}[leftmargin=1.7em]
  \item Open the Start menu, search for \textbf{Edit environment variables for your account}, and open it.
  \item Under \textbf{User variables}, select \textbf{Path}, choose \textbf{Edit}, then choose \textbf{New}. Do not replace or delete the existing entries.
  \item Add \texttt{C:\textbackslash msys64\textbackslash ucrt64\textbackslash bin}. If MSYS2 was installed in a different folder, use that installation folder followed by \texttt{\textbackslash ucrt64\textbackslash bin}.
  \item Confirm each dialog with \textbf{OK}, then close and reopen PowerShell, Command Prompt, or the IDE so it reads the updated path.
\end{enumerate}

\noindent Verify that Windows finds the compiler:

\begin{lstlisting}[style=utpterminal]
where.exe g++
g++ --version
\end{lstlisting}

\noindent The first command should print a location under
\texttt{ucrt64\textbackslash bin}. If it prints nothing, check that the folder
exists and that the new \texttt{Path} entry points to its \texttt{bin}
subfolder. For why UCRT64 uses this compiler directory, see the
\href{https://www.msys2.org/docs/environments/}{\textbf{MSYS2 environment guide}}.

\noindent If your course requires Microsoft's compiler, install Visual Studio
or Build Tools with the
\href{https://learn.microsoft.com/en-us/cpp/build/vscpp-step-0-installation}{\textbf{Desktop development with C++ workload}} and use its Developer Command Prompt.

\medskip
\noindent\textcolor{UTPBlue!45}{\rule{\linewidth}{0.5pt}}\par
\medskip
\noindent\textbf{Linux}\par
\noindent On Ubuntu or Debian, install and check GNU C++. Other distributions
provide GCC/G++ through their package manager:

\begin{lstlisting}[style=utpterminal]
sudo apt update
sudo apt install g++
g++ --version
\end{lstlisting}

\medskip
\noindent\textcolor{UTPBlue!45}{\rule{\linewidth}{0.5pt}}\par
\medskip
\noindent\textbf{macOS}\par
\noindent Install Apple's Command Line Tools (the full Xcode application is
not required), then check that Clang is available:

\begin{lstlisting}[style=utpterminal]
xcode-select --install
clang++ --version
\end{lstlisting}
